\subsection{齐性空间体积的半连续性}

在本 No. 中，对每一测度 \(\alpha \in \Gamma\)，我们置

\[
\mathrm{Q} _ {\alpha} = \mathrm{G} / \mathrm{H} _ {\alpha},\tag{1}
\]

且用 \(\pi_{\alpha}\) 表示典范映射 \(\mathbf{G} \to \mathbf{Q}_{\alpha}\)。

设 \(\Gamma^0\) 为 \(\Gamma\) 中由这样的测度 \(\alpha\) 组成的子集：子群 \(\mathbf{H}_{\alpha}\)（\(\mathbf{G}\) 的）是单模的；\(\Gamma^0\) 的元素由如下事实刻画：\(\alpha(f) = \alpha(\check{f})\) 对每一函数 \(f \in \mathcal{K}(\mathbf{G})\) 成立（\(\mathcal{K}(\mathbf{H}_{\alpha})\) 的每个函数都可由 Urysohn 定理延拓为 \(\mathcal{K}(\mathbf{G})\) 的函数）；由此得出 \(\Gamma^0\) 是 \(\Gamma\) 的一个闭子集。回忆：对每一 \(\alpha \in \Gamma^0\)，商测度 \(\mu_{\alpha} = \mu / \alpha\)（\(\mathbf{Q}_{\alpha}\) 上的）有定义且在 \(\mathbf{G}\) 下相对不变（Ch. VII, §2, No.~6 定理 3）；再回忆：对每一函数 \(f \in \mathcal{K}(\mathbf{G})\)，

\[
\int_ {\mathbf {G}} f (x) d \mu (x) = \int_ {\mathbf {Q} _ {\alpha}} d \mu_ {\alpha} (\dot {x}) \int_ {\mathrm{H} _ {\alpha}} f (x s) d \alpha (s),\tag{2}
\]

其中 \(\dot{x} = \pi_{\alpha}(x)\) 是 \(x\in \mathbf{G}\) 在 \(\mathbf{Q}_{\alpha}\) 中的典范像。

命题 4. --- 设 \(\Gamma^0\) 为测度 \(\alpha \in \Gamma\)（使得 \(\mathrm{H}_{\alpha}\) 是单模的）组成的集合，且对每一 \(\alpha \in \Gamma^0\) 置 \(\mu_{\alpha} = \mu / \alpha\)；则映射 \(\alpha \mapsto \| \mu_{\alpha} \|\)（从 \(\Gamma^0\) 到 \(\overline{\mathbf{R}}\) 中）对淡拓扑是下半连续的。

对每一 \(\alpha \in \Gamma^0\) 与每一函数 \(f\in \mathcal{K}(\mathbf{G})\)，置

\[
f _ {\alpha} (\dot {x}) = \int_ {\mathrm{H} _ {\alpha}} f (x s) d \alpha (s) = (f * \alpha) (x),
\]

其中卷积是相对于右 Haar 测度 \(\mu\) 取的，并且用到了 \(\check{\alpha} = \alpha\) 这一事实（§4, No.~4 公式 (11)）。我们知道（Ch. VII, §2, No.~1 命题 2），映射 \(f \mapsto f_{\alpha}\)（从 \(\mathcal{H}_{+}(G)\) 到 \(\mathcal{H}_{+}(Q_{\alpha})\) 中）是满射；因此，由 (2)，

\[
\| \mu_ {\alpha} \| = \sup _ {f \in \mathcal {K} _ {+} (\mathrm{G}), f \neq 0} \mu_ {\alpha} (f _ {\alpha}) / \| f _ {\alpha} \| = \sup _ {f \in \mathcal {K} _ {+} (\mathrm{G}), f \neq 0} \mu (f) / \| f _ {\alpha} \|,
\]

其中已置

\[
\left\| f _ {\alpha} \right\| = \sup _ {\dot {x} \in \mathrm{Q} _ {\alpha}} | f _ {\alpha} (\dot {x}) | = \sup _ {x \in \mathrm{G}} | (f * \alpha) (x) |.\tag{3}
\]

为建立命题，只需证明：给定 \(f \in \mathcal{K}_{+}(G)\)，映射 \(\alpha \mapsto \|f_{\alpha}\|\) 是淡连续的。现在，设 K 为 f 的支集；函数 \(f * \alpha\) 的支集包含在 \(KH_{\alpha}\) 中且在 \(H_{\alpha}\) 下右不变；从而

\[
\left\| f _ {\alpha} \right\| = \sup _ {x \in K} | (f * \alpha) (x) |.
\]

因此结论由如下事实得出：映射 \(\alpha\mapsto f*\alpha\)（从赋有淡拓扑的 \(\mathcal{M}_{+}(\mathrm{G})\) 到赋有紧收敛拓扑的 \(\mathcal{C}(\mathrm{G})\) 中）是连续的（§4, No.~2 备注 1）。

回忆：若对一个测度 \(\alpha \in \Gamma^0\)，\(\| \mu_{\alpha}\|\) 有限，则 G 必定是单模的（Ch. VII, §2, No.~6 定理 3 的推论 3）。

命题 5. --- 设 g 为一个 \(\mu\)-可积的正数值函数，且设 \(\Gamma^{0}(g)\) 为测度 \(\alpha\in\Gamma^{0}\)（使得 \(\int^{*}g(xs)d\alpha(s)\geqslant1\) 对所有 \(x\in G\)）组成的集合。则映射 \(\alpha\mapsto\|\mu_{\alpha}\|\)（从 \(\Gamma^{0}(g)\) 到 \(\overline{\mathbf{R}}\) 中）是淡连续的。

对每一测度 \(\alpha \in \Gamma^0 (\mathbf{G})\)，回忆（Ch. VII, §2, No.~3 命题 5）函数

\[
g _ {\alpha} (\dot {x}) = \int_ {\mathrm{H} _ {\alpha}} g (x s) d \alpha (s)
\]

\(\mu_{\alpha}\)-几乎处处（在 \(Q_{\alpha}\) 上）有定义，是 \(\mu_{\alpha}\)-可积的，且

\[
\int_ {G} g (x) d \mu (x) = \int_ {Q _ {\alpha}} g _ {\alpha} (\dot {x}) d \mu_ {\alpha} (\dot {x}).\tag{4}
\]

鉴于命题 4，只需证明：在 \(\Gamma^0(g)\) 中，\(\alpha \mapsto \| \mu_\alpha \|\) 是上半连续的。固定一个测度 \(\alpha \in \Gamma^0(g)\)，且设 \(K\) 为 \(G\) 的一个紧子集。在 \(Q_\alpha\) 上存在一个具有紧支集的连续函数，取值于 {[}0,1{]}，在紧集 \(\pi_\alpha(K)\) 上等于 1；由于映射 \(f \mapsto f_\alpha\)（从 \(\mathcal{H}_+(\mathrm{G})\) 到 \(\mathcal{H}_+(\mathrm{Q}_\alpha)\) 中）是满射（Ch. VII, §2, No.~1 命题 2），可见存在一个函数 \(f \in \mathcal{H}_+(\mathrm{G})\)，使得

\[
(f * \alpha) (x) = \int_ {\mathrm{G}} f (x s)   d \alpha (s) \left\{ \begin{array}{l} \leqslant 1 \text {for all} x \in \mathrm{G} \\ = 1 \text {for all} x \in \mathrm{K}. \end{array} \right.
\]

由于 \(\beta \mapsto f * \beta\) 是一个连续映射（从赋有淡拓扑的 \(\mathcal{M}_{+}(\mathbf{G})\) 到赋有紧收敛拓扑的 \(\mathcal{C}(\mathbf{G})\) 中）（§4, No.~2 备注 1），可见对每一 \(\varepsilon > 0\)，\(\mathrm{U}_{\varepsilon}\)（由 \(\beta \in \Gamma^0 (\mathbf{G})\) 中使得

\[
f _ {\beta} (\dot {x}) = \int_ {\mathrm{G}} f (x s) d \beta (s) > 1 - \varepsilon \quad \text { for   all } x \in \mathrm{K}
\]

的测度组成）是 \(\alpha\) 在 \(\Gamma^0 (g)\) 中的一个开邻域；于是对每一 \(\beta \in \mathrm{U}_{\varepsilon}\)，由公式 (2)，我们有

\[
\left\| \mu_ {\alpha} \right\| \geqslant \int_ {G} f (x) d \mu (x) = \int_ {Q _ {\beta}} f _ {\beta} (\dot {x}) d \mu_ {\beta} (\dot {x}) \geqslant (1 - \varepsilon) \mu_ {\beta} \left(\pi_ {\beta} (K)\right).\tag{5}
\]

给定一个数 \(\varepsilon > 0\)，选取一个函数 \(h \in \mathcal{H}_{+}(\mathrm{G})\)，使得 \(\int_{\mathrm{G}} |g(x) - h(x)| d\mu(x) \leqslant \varepsilon\)，并在上文中取 \(\mathrm{K} = \operatorname{Supp}(h)\)。对每一 \(\beta \in \Gamma^0(g)\)，按假设 \(g_\beta(\dot{x}) \geqslant 1\)（对 \(\mu_\beta\)）在 \(\mathbf{Q}_\beta\) 中几乎处处成立，因此

\[
\mu_ {\beta} \big (\mathrm{Q} _ {\beta} - \pi_ {\beta} (\mathrm{K}) \big) \leqslant \int_ {\mathrm{Q} _ {\beta} - \pi_ {\beta} (\mathrm{K})} g _ {\beta} (\dot {x}) d \mu_ {\beta} (\dot {x}) = \int_ {\mathrm{G} - \mathrm{KH} _ {\beta}} g (x) d \mu (x)
\]

由 (4)；由于 \(h\) 在 \(\mathbf{K}\) 之外为零，更不用说在 \(\mathrm{KH}_{\beta}\) 之外为零，由此得出

\[
\begin{array}{c} \mu_ {\beta} \big (\mathrm{Q} _ {\beta} - \pi_ {\beta} (\mathrm{K}) \big) \leqslant \int_ {\mathrm{G} - \mathrm{KH} _ {\beta}} | g (x) - h (x) | d \mu (x) \\ \leqslant \int_ {\mathrm{G}} | g (x) - h (x) | d \mu (x) \leqslant \varepsilon ; \end{array}
\]

把这一结果与 (5) 结合起来，可见

\[
\| \mu_ {\beta} \| \leqslant \varepsilon + \| \mu_ {\alpha} \| / (1 - \varepsilon)
\]

当 \(\beta \in \mathrm{U}_{\varepsilon}\) 时，这就完成了证明。

推论 1. --- 设 K 为 G 的一个紧子集，V 为 e 在 G 中的一个对称紧邻域，c 为一个实数 \textgreater{} 0。映射 \(\alpha \mapsto \|\mu_{\alpha}\|\) 在 \(\alpha \in \Gamma^0\)（满足 \(G = KH_{\alpha}\) 且 \(\alpha(V) \geqslant c\)）组成的集合上的限制是淡连续的。

因为，设 \(g \in \mathcal{K}_{+}(\mathbf{G})\) 为一个函数，使得 \(g(x) \geqslant 1 / c\)（对 \(x \in \mathrm{KV}\)）。对每一 \(x \in \mathbf{K}\)，

\[
\int g (x s) d \alpha (s) \geqslant \int_ {\mathrm{V}} g (x s) d \alpha (s) \geqslant 1
\]

对满足陈述中条件的 \(\alpha\)；此外，由于 \(\pi_{\alpha}(\mathbf{K}) = \mathrm{Q}_{\alpha}\)，有 \(\alpha \in \Gamma^0(g)\)，由此得出推论。

推论 2. --- 设 A 为 G 的一个 \(\mu\)-可积子集。映射 \(\alpha \mapsto \|\mu_{\alpha}\|\) 在集合 \(N_A\)——G 的离散子群 H 的规范化 Haar 测度中使得 \(G = AH\) 的那些——上的限制是淡连续的。

对 \(a \in \mathbf{A}\) 与 \(\alpha \in \mathbf{N}_{\mathbf{A}}\)，

\[
\int \varphi_ {\mathrm{A}} (a s) d \alpha (s) \geqslant \varphi_ {\mathrm{A}} (a) = 1,
\]

而由于 \(\pi_{\alpha}(\mathbf{A}) = \mathbf{Q}_{\alpha}\)，有 \(\mathrm{N}_{\mathrm{A}} \subset \Gamma^{0}(\varphi_{\mathrm{A}})\)，因此推论由命题 5 得出。

